% Part: methods
% Chapter: proofs
% Section: proof-by-contradiction

\documentclass[../../../include/open-logic-section]{subfiles}

\begin{document}

\olfileid[pt]{mth}{prf}{con}

\olsection{Prova por Contradição} 

Em primeira instância, a prova por contradição é um padrão de inferência
usado para provar afirmações negativas.  Suponha que você queira mostrar que
alguma afirmação~$p$ é \emph{falsa}, isto é, você quer mostrar~$\lnot p$.  A
estratégia mais promissora é (a) supor que $p$~é verdadeira, e (b)
mostrar que essa suposição leva a algo que você sabe ser falso.
``Algo que se sabe ser falso'' pode ser um resultado que entra em conflito
com---contradiz---o próprio~$p$, ou alguma outra hipótese da
afirmação geral que você está considerando.  Por exemplo, uma prova de ``se $q$
então $\lnot p$'' envolve supor que $q$~é verdadeira e provar~$\lnot
p$ a partir disso. Se você prova $\lnot p$ por contradição, isso significa
supor $p$ além de~$q$. Se você consegue provar $\lnot q$ a partir de $p$,
mostrou que a suposição~$p$ leva a algo que
contradiz sua outra suposição~$q$, já que $q$~e $\lnot q$ não podem
ser ambas verdadeiras.  É claro que você tem que usar outros padrões de inferência em
sua prova da contradição, bem como desdobrar definições.
Vamos considerar um exemplo.

\begin{prop}
  Se $A \subseteq B$ e $B = \emptyset$, então $A$ não tem !!{element}s.
\end{prop}

\begin{proof}
  Suponha $A \subseteq B$ e $B = \emptyset$. Queremos mostrar que
  $A$ não tem !!{element}s.
  \begin{quote}
    Como esta é uma afirmação condicional, supomos o antecedente e
    queremos provar o consequente. O consequente é: $A$ não tem
    !!{element}s. Podemos tornar isso um pouco mais explícito: não é o
    caso de que haja um~$x \in A$.
  \end{quote}
  $A$ não tem !!{element}s sse não é o caso de que haja um~$x$
  tal que $x \in A$.
  \begin{quote}
    Então determinamos que o que queremos provar é, na verdade, uma
    afirmação negativa~$\lnot p$, a saber: não é o caso de que haja
    um $x \in A$.  Para usar a prova por contradição, temos que supor a
    afirmação positiva correspondente~$p$, isto é, há um $x \in A$, e
    provar uma contradição a partir dela.  Indicamos que estamos fazendo uma
    prova por contradição escrevendo ``por
    absurdo, suponha'' ou mesmo apenas ``suponha que não,'' e então enunciamos
    a suposição~$p$.
  \end{quote}
  Suponha que não: há um $x \in A$.
  \begin{quote}
    Esta é agora a nova suposição que usaremos para obter uma
    contradição. Temos mais duas suposições: que $A \subseteq B$
    e que $B = \emptyset$. A primeira nos dá que $x \in B$:
  \end{quote}
  Como $A \subseteq B$, $x \in B$.
  \begin{quote}
    Mas como $B = \emptyset$, todo !!{element} de $B$ (por exemplo, $x$)
    também deve ser !!a{element} de~$\emptyset$.
  \end{quote}
  Como $B = \emptyset$, $x \in \emptyset$. Isso é uma contradição,
  já que, por definição, $\emptyset$ não tem !!{element}s.
  \begin{quote}
    Isso já completa a prova: chegamos ao que precisávamos (uma
    contradição) a partir das suposições que montamos, e isso significa
    que as suposições não podem ser todas verdadeiras. Como as duas primeiras
    suposições ($A \subseteq B$ e $B = \emptyset$) não são
    contestadas, deve ser a última suposição introduzida (há um
    $x \in A$) que tem que ser falsa. Mas, se quisermos ser minuciosos,
    podemos explicitar isso.
  \end{quote}
  Assim, nossa suposição de que há um $x \in A$ tem que ser falsa, logo,
  $A$ não tem !!{element}s, por prova por contradição.
\end{proof}

Toda afirmação positiva é trivialmente equivalente a uma afirmação negativa: $p$
sse $\lnot\lnot p$.  Assim, as provas por contradição também podem ser usadas para
estabelecer afirmações positivas ``indiretamente,'' como segue: Para provar~$p$,
leia-o como a afirmação negativa $\lnot\lnot p$. Se conseguirmos provar uma
contradição a partir de $\lnot p$, estabelecemos $\lnot\lnot p$ por
prova por contradição, e portanto~$p$.

No último exemplo, buscávamos provar uma afirmação negativa, a saber, que
$A$ não tem !!{element}s, e por isso a suposição que fizemos para os fins
da prova por contradição (isto é, que há um $x \in A$) era uma
afirmação positiva. Ela nos deu algo com que trabalhar, a saber, o
hipotético $x \in A$ sobre o qual continuamos a raciocinar até chegarmos
a $x \in \emptyset$.

Ao provar uma afirmação positiva indiretamente, a suposição que você faria
para os fins da prova por contradição seria negativa. Mas muito
frequentemente você pode facilmente reformular uma afirmação positiva como uma afirmação negativa,
e uma afirmação negativa como uma afirmação positiva.  Nossa prova anterior teria
sido essencialmente a mesma se tivéssemos provado ``$A = \emptyset$''
em vez do consequente negativo ``$A$ não tem !!{element}s.'' (Pela
definição de $=$, ``$A = \emptyset$'' é uma afirmação geral, já que
ela se desdobra em ``todo !!{element} de~$A$ é !!a{element} de~$\emptyset$
e vice-versa''.) Mas é facilmente visto que ela é equivalente à
afirmação negativa ``não: há um $x \in A$.''

Então, às vezes é mais fácil trabalhar com $\lnot p$ como suposição do
que provar~$p$ diretamente.  Mesmo quando uma prova direta é igualmente
simples ou até mais simples (como nos próximos exemplos), algumas pessoas preferem
proceder indiretamente.  Se a dupla negação o confunde, pense em uma
prova por contradição de alguma afirmação como uma prova de uma contradição
a partir da afirmação \emph{oposta}. Assim, uma prova por contradição de $\lnot
p$ é uma prova de uma contradição a partir da suposição~$p$; e a prova por
contradição de~$p$ é uma prova de uma contradição a partir de~$\lnot p$.

\begin{prop}
$A \subseteq A \cup B$.
\end{prop}

\begin{proof}
  Queremos mostrar que $A \subseteq A \cup B$.
  \begin{quote}
    À primeira vista, esta é uma afirmação positiva: todo $x \in A$ está
    também em $A \cup B$.  A negação disso é: algum $x \in A$ está
    $\notin A \cup B$. Então podemos provar a afirmação indiretamente
    supondo essa afirmação negada, e mostrando que ela leva a uma
    contradição.
    \end{quote}
  Suponha que não, isto é, $A \nsubseteq A \cup B$.
  \begin{quote}
    Temos uma definição de $A \subseteq A \cup B$: todo $x \in A$ está
    também $\in A \cup B$.  Para entender o que $A \nsubseteq A \cup B$
    significa, temos que usar alguma manipulação lógica elementar sobre a
    definição desdobrada: é falso que todo $x \in A$ esteja também $\in
    A \cup B$ sse há \emph{algum}~$x \in A$ que está $\notin C$.
    (Este é um ponto em que você quer ser muito cuidadoso: muitas tentativas de
    prova por contradição de estudantes falham porque eles analisam a
    negação de uma afirmação como ``todos os $A$ são $B$'' incorretamente.) Em
    outras palavras, $A \nsubseteq A \cup B$ sse há um $x$ tal que
    $x \in A$ e $x \notin A \cup B$. A partir daí, é fácil.
  \end{quote}
  Então, há um $x \in A$ tal que $x \notin A \cup B$.  Pela
  definição de $\cup$, $x \in A \cup B$ sse $x \in A$ ou $x \in
  B$. Como $x \in A$, temos $x \in A \cup B$. Isso contradiz a
  suposição de que $x \notin A \cup B$. 
\end{proof}

\begin{prob}
Prove \emph{indiretamente} que $A \cap B \subseteq A$.
\end{prob}

\begin{prop}
Se $A \subseteq B$ e $B \subseteq C$ então $A \subseteq C$.
\end{prop}

\begin{proof}
  Suponha $A \subseteq B$ e $B \subseteq C$. Queremos mostrar $A
  \subseteq C$.
  \begin{quote}
    Vamos proceder indiretamente: supomos a negação do que queremos
    estabelecer.
  \end{quote}
  Suponha que não, isto é, $A \nsubseteq C$.
  \begin{quote}
    Como antes, raciocinamos que $A \nsubseteq C$ sse nem todo $x \in A$
    está também $\in C$, isto é, algum $x \in A$ está $\notin C$.  Não se preocupe,
    com prática você não terá mais que pensar muito para desdobrar
    negações como esta.
  \end{quote}
  Em outras palavras, há um~$x$ tal que $x \in A$ e $x \notin C$.
  \begin{quote}
    Agora podemos usar isso para chegar à nossa contradição. É claro que
    teremos que usar as outras duas suposições para fazê-lo.
  \end{quote}
  Como $A \subseteq B$, $x \in B$. Como $B \subseteq C$, $x \in
  C$. Mas isso contradiz $x \notin C$.
\end{proof}

\begin{prop}
Se $A \cup B = A \cap B$ então $A = B$.
\end{prop}

\begin{proof}
  Suponha $A \cup B = A \cap B$. Queremos mostrar que $A = B$.
  \begin{quote}
    O começo é agora rotineiro:
  \end{quote}
  Suponha, por absurdo, que $A \neq B$.
  \begin{quote}
    Nossa suposição para a prova por contradição é que $A \neq
    B$. Como $A = B$ sse $A \subseteq B$ e $B \subseteq A$, obtemos
    que $A \neq B$ sse $A \nsubseteq B$ \emph{ou} $B \nsubseteq
    A$. (Note como é importante ter cuidado ao manipular
    negações!{}) Para provar uma contradição a partir dessa disjunção,
    usamos uma prova por casos e mostramos que, em cada caso, uma contradição
    se segue.
  \end{quote}
  $A \neq B$ sse $A \nsubseteq B$ ou $B \nsubseteq A$. Distinguimos
  casos.
  \begin{quote}
    No primeiro caso, supomos $A \nsubseteq B$, isto é, para algum $x$,
    $x \in A$ mas $\notin B$. $A \cap B$ é definido como aqueles
    !!{element}s que $A$ e $B$ têm em comum, então, se algo
    não está em um deles, não está na interseção. $A \cup B$ é
    $A$ junto com $B$, então qualquer coisa em um ou outro está também na
    união. Isso nos diz que $x \in A \cup B$ mas $x \notin A \cap
    B$, e logo que $A \cap B \neq A \cup B$.
  \end{quote}
  
  Caso 1: $A \nsubseteq B$. Então, para algum $x$, $x \in A$ mas $x \notin
  B$. Como $x \notin B$, então $x \notin A \cap B$. Como $x \in A$,
  $x \in A \cup B$. Então, $A \cap B \neq A \cup B$, contradizendo a
  suposição de que $A \cap B = A \cup B$.

  Caso 2: $B \nsubseteq A$. Então, para algum $y$, $y \in B$ mas $y \notin
  A$. Como antes, temos $y \in A \cup B$ mas $y \notin A \cap B$, e
  assim $A \cap B \neq A \cup B$, novamente contradizendo $A \cap B = A \cup
  B$. 
\end{proof}

\end{document}
